\subsection{Contador Síncrono MOD 24}

\paragraph{Incremento}

A operação de incremento realiza a transição:
\[
D \mapsto (D + 1) \mod 24
\]

Monta-se a tabela verdade para a função de incremento, considerando as 24 combinações válidas de entrada e suas respectivas saídas:

\begin{center}
\begin{tabular}{c|ccccccc}
$D$ & $O_{11}$ & $O_{10}$ & $O_{03}$ & $O_{02}$ & $O_{01}$ & $O_{00}$ & $Carry$\\
\hline
$0$  & 0 & 0 & 0 & 0 & 0 & 1 & 0\\
$1$  & 0 & 0 & 0 & 0 & 1 & 0 & 0\\
$\vdots$ & & & & & & & \\
$22$ & 1 & 0 & 1 & 1 & 1 & 1 & 0\\
$23$ & 0 & 0 & 0 & 0 & 0 & 0 & 1\\
\end{tabular}
\end{center}

Devido a extensáo do número de bits, a tabela é apresentada de forma resumida. A tabela completa pode ser encontrada no código-fonte do projeto no github (\url{https://url.github.com/}).

A partir da tabela verdade, obtêm-se as seguintes expressões booleanas:

\begin{align*}
O_{11} &= D_{10} D_{03} D_{00} \lor D_{11} \overline{D_{01}} \lor D_{11} \overline{D_{00}} \\
O_{10} &= \overline{D_{10}} D_{03} D_{00} \lor D_{10} \overline{D_{03}} \lor D_{10} \overline{D_{00}} \\
O_{03} &= D_{02} D_{01} D_{00} \lor D_{03} \overline{D_{00}} \\
O_{02} &= \overline{D_{11}} \overline{D_{02}} D_{01} D_{00} \lor D_{02} \overline{D_{01}} \lor D_{02} \overline{D_{00}} \\
O_{01} &= \overline{D_{03}} \overline{D_{01}} D_{00} \lor D_{01} \overline{D_{00}} \\
O_{00} &= \overline{D_{00}} \\
Carry &= D_{11} D_{01} D_{00}
\end{align*}

O sinal $Carry$ indica overflow, isto é, a transição de $23$ para $0$.

As Figuras~\ref{fig:incmod24.png} e~\ref{fig:incmod24-red.png} apresentam a implementação do circuito.

\imgbig{incmod24.png}{Circuito lógico para incremento em módulo 24.}
\imgbig{incmod24-red.png}{Implementação em redstone do circuito de incremento em módulo 24.}

\paragraph{Decremento}

A operação de decremento realiza a transição:
\[
D \mapsto (D - 1) \mod 24
\]

A tabela correspondente é:

\begin{center}
\begin{tabular}{c|ccccccc}
$D$ & $O_{11}$ & $O_{10}$ & $O_{03}$ & $O_{02}$ & $O_{01}$ & $O_{00}$ & $Carry$\\
\hline
$0$  & 1 & 0 & 1 & 0 & 1 & 1 & 1\\
$1$  & 0 & 0 & 0 & 0 & 0 & 0 & 0\\
$\vdots$ & & & & & & & \\
$23$ & 1 & 1 & 1 & 1 & 1 & 1 & 0\\
\end{tabular}
\end{center}

As expressões booleanas minimizadas são:

\begin{align*}
Carry &= \overline{D_{11}} \overline{D_{10}} \overline{D_{03}} \overline{D_{02}} \overline{D_{01}} \overline{D_{00}} \\
O_{11} &= \overline{D_{11}} \overline{D_{10}} \overline{D_{03}} \overline{D_{02}} \overline{D_{01}} \overline{D_{00}} \lor D_{11} D_{00} \lor D_{11} D_{01} \\
O_{10} &= D_{10}(D_{00} \lor D_{01} \lor D_{02} \lor D_{03}) \lor D_{11} \overline{D_{01}} \overline{D_{00}} \\
O_{03} &= D_{03} D_{00} \lor D_{10} \overline{D_{03}} \overline{D_{02}} \overline{D_{01}} \overline{D_{00}} \lor D_{11} \overline{D_{01}} \overline{D_{00}} \\
O_{02} &= D_{02} D_{00} \lor D_{02} D_{01} \lor D_{03} \overline{D_{00}} \\
O_{01} &= \overline{D_{11}} \overline{D_{10}} \overline{D_{01}} \overline{D_{00}} \lor D_{01} D_{00} \lor D_{02} \overline{D_{01}} \overline{D_{00}} \lor D_{03} \overline{D_{00}} \\
O_{00} &= \overline{D_{00}}
\end{align*}

Neste caso, o sinal $Carry$ indica underflow, ou seja, a transição de $0$ para $23$.

As Figuras~\ref{fig:decmod24.png} e~\ref{fig:decmod24-red.png} mostram a implementação do circuito.

\imgbig{decmod24.png}{Circuito lógico para decremento em módulo 24.}
\imgbig{decmod24-red.png}{Implementação em redstone do circuito de decremento em módulo 24.}

\paragraph{Considerações de Projeto}

Embora funcional, o contador não é autocorretivo, pois estados inválidos não são redirecionados ao ciclo válido.

Uma abordagem mais robusta consiste em tratar estados inválidos como condições de correção ou como \textit{don't care} na minimização lógica.